\documentclass[aps,prb,onecolumn,nofootinbib,superscriptaddress]{revtex4-2}
\usepackage{amsmath,amssymb,amsthm,mathtools,bm}
\usepackage{booktabs}
\usepackage[colorlinks=true,linkcolor=blue,citecolor=blue,urlcolor=blue]{hyperref}
\usepackage{microtype}

\newcommand{\repo}[1]{\texttt{\detokenize{#1}}}

\begin{document}

\title{Supervisor 05 Reliability and Provenance Synthesis}
\author{Supervisor 05}
\date{June 11, 2026}

\begin{abstract}
This note synthesizes the assigned manifest/data audit and figure-provenance audit, then
adds an independent read-only check of reliability and provenance metadata in the repository.
The robust evidence spine is the structured campaign layer: JSON manifests, run summaries,
run-index CSV files, and analysis manifests that record the canonical dynamics entry point,
sample counts, saved-observable policy, and source-to-derived-product links.  The weakest
layer is the legacy \repo{cache/G_history_samples} archive, which is enormous, mostly
filename-provenanced, and contains rerun, test, and duplicate-style artifacts.  The most
important caveats for the Boss are the postselection sample-count semantics, a
summary-figure mismatch in charge-fluctuation tables, skipped or missing analysis inputs,
hard-coded fallback data in the $\alpha_2$ synthesis panel, and one dirty/empty
\repo{dw_convergence} manifest.
\end{abstract}

\maketitle
\setcounter{tocdepth}{2}
\tableofcontents

\section{Source Scope}

I first read \repo{REPO_POLICY.md}.  I then read only the assigned worker notes:
\repo{repo_synthesis/main_pass/worker_notes/worker_15_manifest_data_audit.tex} and
\repo{repo_synthesis/main_pass/worker_notes/worker_16_figure_table_provenance_caveats.tex}.
I did not read any other worker notes, supervisor notes, or boss notes.  I did not edit
files, execute notebooks, or load large binary arrays.

The independent audit covered text metadata in manifests, run summaries, indices, logs,
CSV tables, \repo{summary_of_results/}, \repo{cpu_data/}, \repo{gpu_data/}, selected
figure-generation scripts, and \repo{cache/} by filename and byte size only.  Every path
containing \repo{haining}, \repo{haoyu}, or \repo{erroneous}, case-insensitive, was
excluded.  The non-excluded scan found 505 manifest or run-summary JSON files; the largest
contributors were \repo{colab_charge_fluctuations} with 301, \repo{colab_small_system_testing}
with 80, \repo{choi_covariance_cpu} with 66, \repo{colab_partial_post-select} with 19, and
\repo{colab_large_entanglement_scaling_N20} with 9.  The non-excluded \repo{cache} layer
contains 98 files, including 38 \repo{.npz}/\repo{.npy} arrays larger than 1 GiB, totaling
about 653 GiB by filename-size metadata.

\section{Synthesis of Worker Claims}

Worker 15's central claim is correct: the newer campaign/run-summary layer is the only
safe primary data spine.  Its strongest records name \repo{classA_U1FGTN.run_markov_circuit}
or \repo{classA_U1FGTN_gpu.run_markov_circuit}, state whether covariance histories were
saved, report requested and actual samples, and preserve relative paths to run products.
Worker 15's warnings about the legacy cache, postselection sample downgrades, and the
dirty \repo{dw_convergence} manifest are directly visible in independent checks.

Worker 16's figure-provenance hierarchy is also correct.  The seven
\repo{summary_of_results/figures} PDFs are regenerated by
\repo{summary_of_results/scripts/make_summary_figures.py}, which declares itself a
synthesis layer and performs no circuit simulation.  CSV-backed summary panels have stronger
provenance than the \repo{trajectory_vs_mean_channel.pdf} panel, which loads two \repo{npz}
files, and much stronger provenance than \repo{alpha2_truncation_limit.pdf}, whose numerical
tables are hard-coded fallbacks copied from an executed notebook state rather than exported
as stable tables.

\section{Independent Repo-Audit Adjustments}

The main adjustment is that the summary purification figure needs sharper qualification.
The plotting script reads scalar purification metrics for both protocols, but the two
charge-statistic analysis manifests feeding the charge-variance and centered-charge panels
select only \repo{N20x30_perfect_correction}, \repo{N20x40_perfect_correction}, and
\repo{N20x50_perfect_correction}.  Thus any plotted or textual postselection
charge-fluctuation comparison in that panel is not supported by the same charge-statistics
tables.

A second adjustment is that partial-postselection data under \repo{colab_partial_post-select}
are well manifested but auxiliary to the summary note.  A representative run summary records
the canonical GPU entry point, no covariance-history saving, Born-rule probability output,
400 scalar rows, and checkpoint manifests.  It is reliable as a separate campaign, but it
should not be silently merged with the perfect-correction/postselection N20 summary claims.

A third adjustment is that not all useful analysis metadata are equally primary.  The
analysis-output tree contains 116 non-excluded CSV/JSON/log text metadata files, including
11 \repo{analysis_manifest.json} files.  These are valuable for derived-product provenance,
but primary claims should still point back to campaign manifests or run summaries whenever
possible.

\section{Key Evidence Table}

\begin{table}[h]
\caption{Representative evidence used in this synthesis.  Line numbers refer to the inspected
repository text files.}
\begin{tabular}{p{0.34\textwidth}p{0.18\textwidth}p{0.40\textwidth}}
\toprule
Source & Lines & Evidence \\
\midrule
\repo{REPO_POLICY.md} & 1--19 & Requires canonical CPU/GPU dynamics entry points and production metadata recording. \\
\repo{summary_of_results/summary_of_results.tex} & 33--47, 383--393 & States the central narrative but explicitly leaves finite-size, correlator, chirality, and $\alpha_2$ table issues open. \\
\repo{summary_of_results/scripts/make_summary_figures.py} & 1--7, 28--71, 592--600 & Single summary figure script; reads explicit products, avoids excluded GPU directory, and generates seven PDFs. \\
\repo{summary_of_results/scripts/make_summary_figures.py} & 210--287 & Purification figure combines scalar metrics with charge-variance and centered-charge CSV tables. \\
\repo{summary_of_results/scripts/make_summary_figures.py} & 460--549 & $\alpha_2$ panel uses fallback tables because notebook output lacks stable CSV export. \\
\repo{colab_charge_fluctuations/gpu_data/purification_dynamics_maxmix/.../campaign_manifest.json} & 11, 92--100, 126--168, 331--363 & Canonical GPU entry point, no covariance-history saving, postselection forced from requested 100 samples to actual 1, and total saved observable bytes. \\
\repo{colab_charge_fluctuations/analysis_outputs/purification_total_charge_variance.../analysis_manifest.json} & 51, 62--70 & Defines charge-variance table but selects only perfect-correction config IDs. \\
\repo{colab_charge_fluctuations/analysis_outputs/purification_total_charge_histogram.../analysis_manifest.json} & 56, 63--73 & Defines centered-charge table but selects only perfect-correction config IDs. \\
\repo{colab_charge_fluctuations/analysis_outputs/streaming_covariance_characterization/.../analysis_manifest.json} & 36--107 & Maps six records to GPU inputs and records skipped N20x60 run due to missing \repo{local_charge_cell.npz} and fit array. \\
\repo{colab_charge_fluctuations/cpu_data/purification_charge_sharpening_alpha_sweep/.../campaign_manifest.json} & 25--35, 47--48 & Strong CPU provenance: canonical CPU entry, complete campaign, no covariance-history saving, active-slab observer, and matched expected/completed samples. \\
\repo{colab_small_system_testing/gpu_data/covariance_protocol_histories/.../campaign_manifest.json} & 8, 217--223, 256--300, 333--377 & Canonical GPU entry, history storage metadata, shard shapes/sizes, and postselection saved as one deterministic sample despite 10 requested samples. \\
\repo{colab_large_entanglement_scaling_N20/.../run_summary.json} & 49--100, 113--172 & Strong large-system run summary: canonical GPU engine, A100 40 GB batching, final-only Markov metadata, fit quality, and complete status. \\
\repo{colab_large_entanglement_scaling_N20/.../fit_rows.csv} & per-file line 2 & Final slopes for $N_y=30,40,50$ are about 0.356, 0.351, and 0.349 with high $R^2$. \\
\repo{colab_small_system_testing/.../pure_state_square_correlations_cpu/log_chord_fit_summary.csv} & 1--13 & Fixed-$x$ slopes are unstable, while $x$-averaged slopes are near 2.1--2.36. \\
\repo{dw_convergence/n20x20_maxmix_dwtrunc_convergence_cpu/analysis_manifest.json} & 2--24, 76--90 & Records canonical CPU configuration, but git status is dirty and \repo{runs} is empty. \\
\repo{cache/G_history_samples} & filename/size only & 38 non-excluded arrays exceed 1 GiB, totaling about 653 GiB; filenames include rerun, testing, extracted, and legacy-duplicate patterns. \\
\bottomrule
\end{tabular}
\end{table}

\section{Contradictions and Caveats}

The largest semantic contradiction is sample counting in postselection.  Several manifests
request ensemble sizes such as 10 or 100, while the canonical GPU postselection driver saves
one deterministic branch.  This is documented rather than hidden, but it must not be averaged
or weighted as if it were an equal-sample ensemble.

The summary purification panel has a provenance mismatch: the text and caption discuss both
perfect correction and postselection charge fluctuations, while the charge-statistic tables
used by the figure script are selected only from perfect-correction configurations.  This
should be scrutinized before any final quantitative claim about postselection charge variance.

The \repo{N20x60} streaming run is explicit negative evidence, not completed evidence: the
analysis manifest records missing \repo{local_charge_cell.npz} and
\repo{streaming_covariance_scaling_fits.npz}.  The current init-default scaling story is
therefore an $N_y=30,40,50$ derived analysis, not a complete N20x60-inclusive series.

The $\alpha_2$ truncation panel is useful for physical interpretation but weak as archival
provenance.  Until its numbers are exported to CSV or JSON, the Boss should describe it as
a diagnostic synthesis panel rather than a primary data table.

The legacy \repo{cache} layer should be treated as archival unless a current manifest points
to a specific file.  Its size, naming heterogeneity, rerun labels, extracted arrays, and
duplicate-style names make it unsuitable as a primary evidence layer by itself.

\section{Confidence}

Confidence is high for metadata-level conclusions: canonical entry-point recording,
postselection sample downgrades, missing-input records, dirty git provenance, figure-script
inputs, and cache-size risk are directly visible from text files and file metadata.

Confidence is moderate for the hierarchy of primary versus derived products.  That judgment
follows from provenance completeness, but it does not numerically validate the arrays.

Confidence is low for claims about numerical reproducibility from raw covariance histories,
because notebooks were not executed and binary arrays were not loaded.

\section{What the Boss Should Preserve or Scrutinize}

Preserve the structured campaign/run-summary layer as the manuscript's data spine, especially
records with canonical engine fields, explicit sample counts, completion status, saved-product
paths, and run-index CSVs.

Preserve the ensemble-theoretic distinction
\[
\overline{O[G_\xi]} \neq O[\overline{G_\xi}] \neq O[G_{\rm ps}],
\]
because both the summary note and the provenance metadata support treating postselection and
mean-channel evolution as controls rather than the physical adaptive ensemble.

Scrutinize any postselection-versus-perfect-correction charge-fluctuation statement, the
$\alpha_2$ hard-wall/truncation panel, correlator exponent precision claims, and any use of
legacy \repo{cache/G_history_samples} files not explicitly tied to a current manifest.

\section{Cross-Supervisor Addendum}

\subsection{Cross-note agreements}

All supervisor notes agree that the manuscript spine should distinguish the physical
trajectory ensemble from controls:
\[
\overline{O[G_\xi]} \neq O[\overline{G_\xi}] \neq O[G_{\rm ps}] .
\]
Perfect correction is Born-sampled adaptive dynamics with deterministic feedback after
mismatches; postselection is a forced branch, often represented by one actual sample.  The
canonical public entry points remain \repo{classA_U1FGTN.run_markov_circuit} and
\repo{classA_U1FGTN_gpu.run_markov_circuit}.

The strongest physics story is hierarchical.  Exact/OW domain-wall target states give the
theory benchmark; trajectory-resolved perfect correction gives the physical preparation
claim; entropy scaling is the cleanest critical diagnostic; correlators, modular chirality,
Lyapunov/Choi gaps, and partial postselection are supporting or diagnostic layers with
specific caveats.

\subsection{Conflicts and corrections to this note}

This provenance note should be sharpened in one respect: a recorded
\repo{canonical_dynamics_entry_point} string proves use of the public method name, but not
byte identity with the current canonical source.  Supervisor 02 finds stale Colab/OSG copies
of \repo{classA_U1FGTN_gpu.py} and a narrower CPU copy drift.  For flagship claims, source
checksum or copied-source identity should accompany the manifest.

The summary charge-fluctuation warning is reinforced, not weakened, by Supervisor 03:
postselection and perfect correction are not equal-sample ensembles in available products.
My earlier caution about the purification charge-statistic tables should therefore be treated
as a manuscript-level quantitative restriction.

Supervisor 04 adds a terminology correction: Lyapunov, Choi-transfer, and Markov-transfer
operators are related diagnostics, not one interchangeable ``gap.''  Provenance tables should
preserve the diagnostic name and endpoint/censoring policy.

\subsection{Boss priorities}

\begin{enumerate}
\item Preserve the trajectory-resolved perfect-correction narrative, with postselection,
mean-channel, Lindblad, Choi, and transfer objects explicitly labeled as controls or distinct
diagnostics.
\item Anchor primary numerical claims to structured campaign/run-summary metadata, but require
source identity for high-stakes results, not only a canonical-entry string.
\item Use entropy scaling as the leading critical evidence; keep correlator exponents,
modular velocities, partial-postselection trends, and transfer-gap scaling caveated.
\item Scrutinize postselection sample counts, summary purification charge panels, GPU Choi
final-cycle claims, CPU Choi endpoint saturation, and the hard-coded \(\alpha_2\) fallback table.
\item Treat legacy \repo{cache/G_history_samples}, stale OSG paths, and dirty/empty manifests
as archival or diagnostic unless a current manifest explicitly resolves their identity.
\end{enumerate}

\end{document}
